% Section 1. Drafted from docs/PAPER-OUTLINE.md section 1, docs/OVERVIEW.md and README.md.

\section{Introduction}\label{sec:intro}

Replicated systems deliver the same events to different replicas in different orders, because networks
delay, reorder and retry messages. Real applications also have rules: an order ships only once it is
paid, a balance stays above a floor, a shelf holds at most so many items. In practice, systems keep such
rules by repairing after the fact. An event is applied, and if it breaks a rule, a deterministic repair
restores it: a saga compensates, a constraint is re-established, a value is capped. The question this
paper answers is when that works: when replicas that have seen the same events and run the same repair
end in the same state, whatever order the events arrive in and whenever the repair runs.

It does not always work. A shelf holds at most 5 items and has 4 in stock; the events ``add 3'' and
``remove 1'' arrive, and the repair caps the stock at 5. One replica applies ``add 3'', caps at 5 and
removes one, ending at 4; another removes one, adds three and caps, ending at 5. The same events and the
same repair give different results.

The two established answers to concurrency sidestep the question. CRDTs~\cite{shapiro2011} design
operations so that order never matters, and need no repair. Invariant
confluence~\cite{bailis2014} asks which operations preserve an invariant without coordination, with the
invariant checked and never repaired. For compensation as practiced, in sagas~\cite{garcia1987} and in
systems that add invariants to weakly consistent replication~\cite{li2012,balegas2015}, we found no
exact condition in that literature for when concurrent orders converge.

\paragraph{The answer in one sentence.} Raw states may differ along the way; what has to agree is the
result after repair. The paper makes this precise at three scales: for a single event, for finite
histories, and for infinite runs.

\paragraph{Contributions.}
\begin{enumerate}[nosep]
  \item Conditions at three scales, with a proved hierarchy between them: state descent, necessary at
    one event; history descent, which with state descent and terminating repair is exact for a single
    registry; and canonical recurrence, exact for infinite runs on finite state spaces
    (\cref{sec:main}).
  \item The relation to CRDTs stated precisely: repair strictly enlarges which operations can be used, but
    every convergent repaired system is a CRDT on its repaired steps; and the relation to invariant
    confluence (\cref{sec:crdt}).
  \item A three-layer framework, for executions, consistent states and long-run behavior, with proved
    interfaces and the counterexamples that mark where each stops, from which the single-registry theorem
    is derived without circularity (\cref{sec:layer-a,sec:layer-c,sec:layer-b,sec:interfaces,sec:base-case}).
  \item Tests of the framework where it could fail: classical theorems from outside the development
    recovered with its files frozen, and two predictions committed before their tests, one a bet about
    matching markets and one a computed prediction about a widely used database, both confirmed
    (\cref{sec:tests}).
  \item A regime map with an exact condition, a hardness result or a listed open problem in every regime,
    all mechanized in Coq with no axioms, and a checker that runs the practical conditions at build
    time (\cref{sec:map,sec:gsm}).
\end{enumerate}

\paragraph{Scope.} ``Exact'' means necessary and sufficient relative to the model each result states;
counterexamples fix the qualifiers they name. The mathematical tools, Newman's lemma, trace theory,
constraint-satisfaction algebra and fair cycles, are classical and credited where they are used
(\cref{sec:related,sec:limits}).
